\section{Convergence of distributed DRS with inexact Gauss--Newton updates}
\label{sec:drs-gn-convergence}

This section separates the exact DRS statement, which follows from
Themelis and Patrinos~\cite{themelis2020douglas}, from the finite
Gauss--Newton (GN) implementation.  A finite number of GN iterations does
not, in general, evaluate a proximal map.  Consequently, convergence of the
implemented method requires an explicit inexactness condition in addition to
the assumptions of exact DRS.

\subsection{Reduced consensus formulation}

Let $c_i$ denote the local copies of the cameras used by cluster $i$, and let
$p_i$ contain the landmarks owned by that cluster.  Define
\begin{equation}
  f_i(c_i,p_i)
  = \frac{1}{2}\sum_{(a,b)\in E_i}
      \|r_{ab}(c_{i,a},p_{i,b})\|^2,
  \qquad
  \bar f_i(c_i) = \inf_{p_i} f_i(c_i,p_i).
  \label{eq:reduced-cluster-objective}
\end{equation}
With $c=(c_1,\ldots,c_K)$ and the camera-consensus subspace
\begin{equation}
  \mathcal C=\{c: c_{i,a}=c_{j,a}
  \text{ whenever camera }a\text{ occurs in clusters }i\text{ and }j\},
\end{equation}
the distributed problem is
\begin{equation}
  \min_c \;\Phi(c):=F(c)+G(c),
  \qquad
  F(c)=\sum_{i=1}^K \bar f_i(c_i),
  \qquad
  G(c)=\delta_{\mathcal C}(c).
  \label{eq:reduced-consensus-problem}
\end{equation}
This reduction is valid when every observation and every landmark is assigned
to exactly one owner.  Duplicating cameras only introduces the constraints in
$\mathcal C$ and does not duplicate the objective.

\paragraph{Standing assumptions.}
The following conditions make the scope of the result explicit.
\begin{enumerate}
  \item[(A1)] The gauge is fixed.  The iterates remain in a compact admissible
  set $\mathcal X$ on which all observed depths are bounded below by a positive
  constant and the camera intrinsics are bounded.
  \item[(A2)] For every $i$ and $c_i\in\mathcal X$, the minimizer
  $p_i^\star(c_i)$ in~\eqref{eq:reduced-cluster-objective} is locally unique,
  remains in a compact set, and satisfies a uniformly nonsingular second-order
  sufficient condition.  Hence $p_i^\star$ is continuously differentiable and
  $\bar f_i$ has an $L_i$-Lipschitz gradient on $\mathcal X$.
  \item[(A3)] There is a globally $C^{1,1}$ function $F_{\rm ext}$ that agrees
  with $F$ on a neighborhood of every point at which the algorithm evaluates
  $F$, its gradient, or its proximal map.  The exact updates for
  $F_{\rm ext}+G$ therefore coincide with those for~\eqref{eq:reduced-consensus-problem}.
  The extended objective is proper, lower semicontinuous, level bounded, and
  attains its minimum.
  \item[(A4)] For the full-sequence conclusion, $F$ and $G$ are
  semialgebraic, or more generally the Lyapunov function used below has the KL
  property.  This holds for the rational BAL projection model restricted away
  from zero depth, squared residuals, polynomial box/gauge constraints, and
  the linear consensus set.  A non-semialgebraic robust loss requires a
  separate definability or KL argument.
\end{enumerate}
Assumption (A2) is the substantive price of eliminating the landmarks.  If it
cannot be justified, the theorem should instead be stated for the unreduced
variables, or smoothness of the reduced functions should be assumed directly.
Assumption (A3) is also stronger than boundedness of the observed iterates:
Themelis and Patrinos assume a globally Lipschitz gradient.  A purely local BA
argument without such an extension requires a localized DRS theorem and
cannot be obtained by citing their Theorem~4.4 alone.

\subsection{Exact fixed-metric DRS}

Let $P\succ0$ be fixed and write $\|x\|_P^2=x^\top P x$.  The metric proximal
map is
\begin{equation}
  \operatorname{prox}^{P}_{h}(s)
  \in\arg\min_x\left\{h(x)+\frac12\|x-s\|_P^2\right\}.
\end{equation}
The exact relaxed DRS iteration is
\begin{align}
  u^k &\in \operatorname{prox}^{P}_{F}(s^k),
  &v^k &\in \operatorname{prox}^{P}_{G}(2u^k-s^k),
  &s^{k+1} &=s^k+\lambda(v^k-u^k).
  \label{eq:metric-drs}
\end{align}

\begin{theorem}[Exact distributed DRS]
\label{thm:exact-distributed-drs}
Suppose (A1)--(A3) hold, $P$ is fixed, and the stepsize and relaxation satisfy
the range in Themelis and Patrinos~\cite[Theorem~4.1]{themelis2020douglas}.
Then the DRS envelope decreases sufficiently, the residuals are square
summable,
\begin{equation}
  \sum_{k=0}^{\infty}\|u^k-v^k\|_P^2<\infty,
  \qquad
  \min_{0\le j\le k}\|u^j-v^j\|_P=o(k^{-1/2}),
  \label{eq:exact-residual-rate}
\end{equation}
and every accumulation point of $(u^k)$ and $(v^k)$ is stationary for
$\Phi$.  If (A4) also holds, both sequences converge to the same stationary
point of~\eqref{eq:reduced-consensus-problem}.
\end{theorem}

\begin{proof}
Set $y=P^{1/2}c$, $\widetilde F(y)=F_{\rm ext}(P^{-1/2}y)$, and
$\widetilde G(y)=G(P^{-1/2}y)$.  Then
\begin{equation}
  P^{1/2}\operatorname{prox}^{P}_{h}(s)
  =\operatorname{prox}_{h\circ P^{-1/2}}(P^{1/2}s),
\end{equation}
so~\eqref{eq:metric-drs} is Euclidean DRS applied to
$\widetilde F+\widetilde G$.  Assumptions (A1)--(A3) imply Assumption I of
Themelis and Patrinos: $\widetilde F$ is continuously differentiable with a
Lipschitz gradient, $\widetilde G$ is proper and lower semicontinuous, and a
minimizer exists.  Their Theorem~4.3 gives sufficient decrease,
square-summability, the best-residual rate in~\eqref{eq:exact-residual-rate},
and stationarity of every accumulation point.  Level boundedness gives
boundedness of the iterates.  An invertible linear coordinate change preserves
semialgebraicity; therefore (A4) permits direct application of their
Theorem~4.4, which proves convergence of the entire primal sequence to one
stationary point.  This concludes the exact version of the proof in
Section~4.4.
\end{proof}

For the classical nonconvex choice $\lambda=1$, a simple conservative
condition is $\gamma<1/(2L)$ in Euclidean coordinates.  Equivalently, for a
scalar metric $P=\rho I$, one may take $\rho>2L$.  For a block metric, the
condition must be imposed after the coordinate change; replacing it by the
componentwise inequalities $\rho_i>2L_i$ is valid only when the transformed
smoothness bound is indeed block separable.

\subsection{Finite Gauss--Newton local solves}

Let $\widehat u^k$ be the output of the finite local GN solve.  Its proximal
optimality error is
\begin{equation}
  e^k=\nabla F(\widehat u^k)+P(\widehat u^k-s^k).
  \label{eq:proximal-optimality-error}
\end{equation}
For the unreduced problem, the corresponding residual includes the landmark
gradient and is evaluated before landmark elimination.  This is preferable in
the implementation because it does not require differentiating
$p_i^\star(c_i)$.

\begin{assumption}[Relative inexactness]
\label{ass:relative-inexactness}
There is an $\eta\ge0$ such that every accepted local update satisfies
\begin{equation}
  \|e^k\|_{P^{-1}}
  \le \eta\,\|\widehat u^k-\widehat v^k\|_P,
  \label{eq:relative-inexactness}
\end{equation}
where $\widehat v^k=\operatorname{prox}^{P}_{G}(2\widehat u^k-s^k)$.
Alternatively, it is sufficient to require
$\sum_k\|e^k\|_{P^{-1}}<\infty$.
\end{assumption}

\begin{lemma}[Perturbed envelope decrease]
\label{lem:perturbed-envelope}
Let $a>0$ be the sufficient-decrease constant of exact DRS in the squared
$P$-norm.  Suppose the inexact update obeys
\begin{equation}
  \mathcal E_P(s^{k+1})
  \le \mathcal E_P(s^k)
     -a\|\widehat u^k-\widehat v^k\|_P^2
     +b\|e^k\|_{P^{-1}}\,
       \|\widehat u^k-\widehat v^k\|_P
     +d\|e^k\|_{P^{-1}}^2
  \label{eq:perturbed-envelope-bound}
\end{equation}
for constants $b,d\ge0$.  If Assumption~\ref{ass:relative-inexactness} holds
with
\begin{equation}
  \widehat a:=a-b\eta-d\eta^2>0,
  \label{eq:accuracy-threshold}
\end{equation}
then
\begin{equation}
  \mathcal E_P(s^{k+1})
  \le \mathcal E_P(s^k)
     -\widehat a\|\widehat u^k-\widehat v^k\|_P^2.
  \label{eq:certified-inexact-decrease}
\end{equation}
\end{lemma}

\begin{proof}
Insert~\eqref{eq:relative-inexactness} into
\eqref{eq:perturbed-envelope-bound} and collect the coefficient of
$\|\widehat u^k-\widehat v^k\|_P^2$.
\end{proof}

The role of Proposition~9 in the original analysis is now precise: its GN
model bounds should be used to establish
\eqref{eq:perturbed-envelope-bound} and to give a computable lower bound on
$\rho_i$ for which~\eqref{eq:accuracy-threshold} holds.  A decrease bound for
the augmented Lagrangian alone does not prove
\eqref{eq:relative-inexactness}, nor does it make a finite GN update an exact
proximal point.

\begin{theorem}[Conditional convergence with finite GN]
\label{thm:inexact-drs-gn}
Suppose (A1)--(A3), Assumption~\ref{ass:relative-inexactness}, and
Lemma~\ref{lem:perturbed-envelope} hold for a fixed $P$.  Suppose also that
the stationarity residual is bounded by
\begin{equation}
  \operatorname{dist}(0,\widehat\partial\Phi(\widehat v^k))
  \le q\|\widehat u^k-\widehat v^k\|_P
      +r\|e^k\|_{P^{-1}}
  \label{eq:stationarity-residual-bound}
\end{equation}
for constants $q,r$.  Then
\begin{equation}
  \sum_{k=0}^{\infty}
    \|\widehat u^k-\widehat v^k\|_P^2<\infty,
  \qquad
  e^k\to0,
\end{equation}
and every accumulation point is stationary for $\Phi$.  If (A4) holds and the
standard KL relative-error bound is satisfied, the entire primal sequence
converges to one stationary point.
\end{theorem}

\begin{proof}
The envelope is lower bounded under (A1)--(A3).  Telescoping
\eqref{eq:certified-inexact-decrease} proves square-summability of the DRS
residual, hence its convergence to zero.  Relative inexactness gives
$e^k\to0$, and~\eqref{eq:stationarity-residual-bound} makes the limiting
stationarity residual vanish.  Boundedness supplies accumulation points and
outer semicontinuity of the limiting subdifferential proves their
stationarity.  Under (A4), the usual uniformized KL argument applied to the
perturbed envelope and its relative-error bound gives finite length and
full-sequence convergence.
\end{proof}

The final KL sentence is a new inexact-DRS argument; it is not a direct
application of Theorem~4.4 of Themelis and Patrinos.  Theorem~4.4 directly
concludes only Theorem~\ref{thm:exact-distributed-drs}.  This distinction
should remain explicit in the paper.

\subsection{Safeguarded acceleration and changing metrics}

An accelerated state $s_{\rm acc}^{k+1}$ is treated as a proposal.  Accept it
only when the same metric, objective, and proximal residual used by the local
step certify
\begin{equation}
  \mathcal E_P(s_{\rm acc}^{k+1})
  \le \mathcal E_P(s^k)
     -\widehat a\|\widehat u^k-\widehat v^k\|_P^2.
  \label{eq:acceleration-safeguard}
\end{equation}
Otherwise execute the certified unaccelerated update.  Since every accepted
or fallback step obeys the same decrease inequality, the telescoping and
stationarity proof of Theorem~\ref{thm:inexact-drs-gn} is unchanged.  This
covers Nesterov, Anderson, or quasi-Newton proposals without claiming that an
unsafeguarded accelerated map is itself DRS.

The preceding theorems assume a fixed metric.  They apply unchanged when a
fixed full block metric is viewed as a linear coordinate transformation.  For
iteration-dependent metrics $P_k$, one additionally needs uniform spectral
bounds
\begin{equation}
  0<mI\preceq P_k\preceq MI
\end{equation}
and a controlled variation condition such as
\begin{equation}
  P_{k+1}\preceq(1+\xi_k)P_k,
  \qquad \sum_k\xi_k<\infty.
\end{equation}
Without such a condition, Theorems~\ref{thm:exact-distributed-drs} and
\ref{thm:inexact-drs-gn} justify the frozen-metric variant, but not arbitrary
per-iteration updates of the $9\times9$ camera blocks.

\paragraph{Quantities to report.}
The convergence claim can be audited experimentally by logging
$\|s^{k+1}-s^k\|_P$, $\|\widehat u^k-\widehat v^k\|_P$,
$\|e^k\|_{P^{-1}}$, their ratio in~\eqref{eq:relative-inexactness}, the merit
decrease in~\eqref{eq:certified-inexact-decrease}, and the metric eigenvalue
and variation bounds.  Consensus error alone is not a fixed-point or
stationarity certificate.